\Needspace{5in}
\begin{prob}
    Below is the \python{merge} function from the \python{mergesort} algorithm,
    with an added \python{print} statement:

    \begin{minted}[autogobble]{python}
        def merge(left, right, out):
            """Merge sorted arrays, store in out."""
            print(len(left))
            print(len(right))
            left.append(float('inf'))
            right.append(float('inf'))
            left_ix = 0
            right_ix = 0

            for ix in range(len(out)):
                if left[left_ix] < right[right_ix]:
                    out[ix] = left[left_ix]
                    left_ix += 1
                else:
                    out[ix] = right[right_ix]
                    right_ix += 1
    \end{minted}

    Suppose this version of \python{merge} is used in \python{mergesort}, and 
    \python{mergesort} is called on the following list:

    \begin{center} 
        \python{[6, 2, 7, 1, 9, 3, 8, 5]}
    \end{center}

    If all of the numbers that are printed are added up, what will the sum be?

    \inlineresponsebox[.75in]{24}
\end{prob}
